
\subsubsection*{2.4.2}
\phantomsection
\label{sga3:VIIB:2.4.2}

Retain the preceding notation and suppose that \(H\) and \(G\) are
topologically flat over \(k\). Then \(G\) is topologically flat both
over \(k\) and over \(G/H\); hence, by \SGARefLink{sga3:VIIB:1.3.3}{1.3.3}, \(G/H\) is
topologically flat over \(k\). Consequently, by \SGARefLink{sga3:VIIB:2.4}{2.4}, the canonical
projection
\[
  q:K\longrightarrow\operatorname{Coker}(f')
\]
is topologically flat, and \(f'\) is an isomorphism from \(G/H\) onto
\(\ker(q)\). It is also clear that
\[
  \operatorname{Coker}(f)
  =
  \operatorname{Coker}(f').
\]
Thus, under the hypothesis that \(G\) and \(H=\ker(f)\) are
topologically flat over \(k\), we have obtained an isomorphism between
\(\ker(q)\), the image of \(f\), and
\[
  G/\ker(f),
\]
the coimage of \(f\). This proves the following theorem.

\medskip
\noindent\textbf{Theorem.}
\phantomsection
\label{sga3:VIIB:2.4.2:theorem}
{\itshape Let \(k\) be a field. The commutative formal \(k\)-groups
form an abelian category.}

\medskip
\noindent\textbf{Corollary.}
\phantomsection
\label{sga3:VIIB:2.4.2:corollary}
{\itshape Let \(k\) be a field. The commutative affine
\(k\)-group schemes form an abelian category.}

This follows from the theorem and the equivalence of categories in
\SGARefLink{sga3:VIIB:2.2.2}{2.2.2}.
